
% TÍTULO CENTRADO
\setcounter{chapter}{1} 
{\centering
{\large\bfseries CAPÍTULO 1}\par\vspace{0.3cm}
{\Large\bfseries INTRODUCCIÓN}\par}
\vspace{1cm}
\addcontentsline{toc}{chapter}{1. CAPÍTULO 1: INTRODUCCIÓN}

% PÁRRAFO DE INTRODUCCIÓN
Pendiente de momento

\section{Antecedentes}
La programación orientada a objetos (POO) es considerada, uno de los 
paradigmas de desarrollo de software moderno más influyentes y extendidos en 
la actualidad. Ventajas competitivas como la reutilización de código, 
la modularidad y la escalabilidad la posicionan como una de las competencias 
más requeridas en el mercado laboral. Por lo tanto, en disciplinas afines a 
la informática y la ingeniería de software, resulta imperativo que los 
estudiantes dominen sus conceptos abstractos. En este contexto, tecnologías 
como la Realidad Aumentada (RA) ofrecen nuevas posibilidades para representar 
visualmente conceptos complejos mediante experiencias interactivas e 
inmersivas.

El uso de la RA para la enseñanza de las ciencias de la computación ha sido 
objeto de estudio en los últimos años. Al respecto, \textcite{Rizwan2019} 
desarrollaron una aplicación móvil basada en RA para visualizar los 
conceptos de POO y comprobar sus interacciones dinámicas. Su investigación 
concluyó que, a pesar de ser una tecnología aún novedosa, la RA genera 
resultados óptimos al facilitar la transferencia de conocimiento abstracto 
de una manera más práctica a los estudiantes.

Por su parte, \textcite{Abidin2020} diseñaron, desarrollaron y 
evaluaron OOP-AR, una aplicación AR para que los estudiantes puedan aprender 
POO de manera atractiva. Tras una fase de pruebas con una muestra de 20 
participantes, más del 60 \% manifestó un alto grado de conformidad con la 
herramienta. Los hallazgos demostraron que este tipo de tecnologías son 
necesarias para profundizar la comprensión conceptual de POO.

Asimismo, \textcite{kar2022} presentaron el proyecto ORIENT, una 
experiencia de RA intuitiva diseñada para mitigar la frustración académica 
mediante la entrega de contenidos fragmentados y un ciclo de retroalimentación 
positivo. Los resultados experimentales confirmaron que la aplicación posee 
un elevado potencial de transferencia al aula, en cursos introductorios de 
informática.

Finalmente, \textcite{Hang2023} propusieron AR-OOP, una aplicación de 
aprendizaje móvil basado en RA, para ayudar a estudiantes principiantes a 
aprender conceptos básicos de POO. La evaluación del sistema, aplicada a una 
muestra de 30 participantes, arrojó una tasa de aceptación positiva superior 
e igual al 90 \% en términos de usabilidad, interfaz de usuario, satisfacción 
de uso y rendimiento.   

\section{Descripción del problema}
La programación orientada a objetos (POO) es un paradigma que se imparte 
durante los primeros años en carreras afines a las ciencias de la computación. 
Sin embargo, aprenderlo requiere un cambio profundo en la forma de pensar. 
Al respecto, \textcite{Batiha2023} señalan que los estudiantes suelen 
experimentar una alta carga cognitiva ante contenidos de la POO, percibiendo 
la asignatura como compleja y de difícil asimilación. En este sentido, 
\textcite{Shmallo2025} identifican en sus estudios preliminares que los 
estudiantes presentan múltiples errores y nociones equívocas sobre las bases 
de este paradigma. Asimismo, los resultados de \textcite{Abbasi2021} confirman 
que estas dificultades y conceptos erróneos persisten en la comprensión de 
prácticamente todos sus componentes fundamentales. 

El problema de la asimilación de conceptos de POO en estudiantes también se 
puede ver reflejado en estudios realizados en Latinoamérica. 
\textcite{Quirino2024} concluyen que un porcentaje superior al 50 \% de su 
muestra no logra entender el paradigma y sus pilares esenciales, lo que 
repercute en altas tasas de reprobación.

En Bolivia, de manera similar, los estudiantes de las carreras de Ingeniería 
de Sistemas e Ingeniería Informática de la Universidad Mayor de San Simón 
(UMSS) presentan dificultades en la asimilación de conceptos de POO. Este 
problema se traduce en la complejidad que experimenta el alumnado para 
comprender los conceptos abstractos del paradigma e integrarlos con sus 
conocimientos previos. Nociones fundamentales referentes a clases, objetos, 
abstracción, encapsulamiento, polimorfismo y herencia son los que generan un 
mayor índice de incertidumbre al momento de ser implementados de forma 
práctica.

Por otra parte, factores como la baja comprensión lectora, el déficit en el 
razonamiento lógico-matemático, los métodos de enseñanza de los diferentes 
docentes e incluso el propio desinterés del estudiante figuran como algunas 
de las causas más comunes a esta problemática, desencadenando un bajo 
rendimiento académico generalizado, altas tasas de reprobación y abandono, 
así como un estado de desmotivación en el alumnado.


\subsection{Definición del problema}
El problema central se define como la baja asimilación de los conceptos 
fundamentales de la Programación Orientada a Objetos en la asignatura de 
Introducción a la Programación. Esta dificultad, atribuida a la alta carga 
cognitiva del paradigma genera nociones conceptuales erróneas que 
desencadenan un bajo rendimiento académico y elevadas tasas de reprobación y 
abandono.

\section{Pregunta de investigación}
¿Cuál es la influencia de la Realidad Aumentada móvil con marcador y sin 
marcador como apoyo en la asimilación de conceptos de Programación Orientada 
a Objetos?

\section{Hipótesis}
La Realidad Aumentada móvil con marcador y sin marcador influye de manera 
positiva como apoyo en la asimilación de conceptos de Programación Orientada 
a Objetos, donde el modelo con marcador facilita la comprensión de estructuras 
fijas como clases y atributos, mientras que el modelo sin marcador mejora el 
entendimiento de procesos dinámicos como la instanciación e interacción entre 
objetos.

\section{Objetivos}
A continuación se presentan el objetivo general y los objetivos específicos.

\subsection{Objetivo general}
Desarrollar dos prototipos funcionales de Realidad Aumentada con marcador y 
sin marcador para determinar su influencia como apoyo en la asimilación de 
conceptos de Programación Orientada a Objetos.

\subsection{Objetivos específicos}
\begin{enumerate}
    \item Identificar las dificultades de los estudiantes en el aprendizaje de
          Programación Orientada a Objetos.
    \item Identificar los conceptos fundamentales de Programación Orientada a 
          Objetos que se puedan abordar en entornos visuales.
    \item Elaborar un prototipo funcional de Realidad Aumentada móvil sin 
          marcador, dirigido a la asimilación de los conceptos fundamentales de 
          Programación Orientada a Objetos identificados.
    \item Elaborar un prototipo funcional de Realidad Aumentada móvil con 
          marcador, dirigido a la asimilación de los conceptos fundamentales de  
          Programación Orientada a Objetos identificados.
    \item Evaluar la influencia de cada modelo de Realidad Aumentada móvil para 
          la asimilación de conceptos fundamentales de  Programación Orientada a 
          Objetos.
    \item Determinar los factores que influyen en cada modelo de Realidad 
          Aumentada móvil para la asimilación de conceptos fundamentales de 
          Programación Orientada a Objetos, a partir de los resultados obtenidos.
\end{enumerate}

\section{Justificación}
El presente trabajo de grado surge a partir de la necesidad de fortalecer la 
comprensión de los conceptos de la POO e identificar los factores que 
influyen en la asimilación de sus fundamentos, con el fin de evaluar en qué 
medida esta se ve favorecida por el uso de la RA como herramienta interactiva. 
La decisión de desarrollar dos prototipos funcionales, uno basado en el 
seguimiento mediante marcadores y otro en el posicionamiento espacial sin 
marcador, responde al propósito de aislar y evaluar el impacto pedagógico de 
cada modelo. Puesto que ambas modalidades implican arquitecturas de 
desarrollo y patrones de interacción notablemente distintos, el proyecto 
requiere una doble fase de modelado 3D, optimización de motores y pruebas de 
usabilidad. Así, la complejidad técnica y el rigor metodológico requeridos 
para calibrar y evaluar simultáneamente dos soluciones independientes 
justifican un esfuerzo colaborativo e interdependiente.

Se escogió la Programación Orientada a Objetos como foco de estudio porque 
constituye uno de los paradigmas fundamentales en el desarrollo de software 
\parencite{Onu2024}. De hecho, esta industria ha experimentado una transición 
masiva desde el estilo procedimental hacia el enfoque orientado a objetos 
\parencite{Agu2022}, lo que respalda los hallazgos de \textcite{Surakka2007} 
sobre su creciente importancia. En este sentido, \textcite{Cummings2020} 
identificaron a la POO como una de las competencias más relevantes y de 
mayor demanda dentro de las 20 habilidades principales de un desarrollador 
de software, de modo que una mejor comprensión de este paradigma se traduce 
en profesionales más capacitados para integrarse a la industria. Sin embargo, 
\textcite{kar2022} sostienen que la POO implica un nivel de abstracción mayor 
en comparación con otros paradigmas, siendo más exigente en los procesos de 
análisis y diseño conceptual.

La literatura hace hincapié en el uso de herramientas didácticas para mejorar 
la comprensión de conceptos en diversas disciplinas. Por esta razón, se 
seleccionó a la Realidad Aumentada como tecnología de apoyo para la 
asimilación de conceptos de POO, debido a que autores como 
\textcite{Batiha2023} señalan que los entornos interactivos contribuyen a 
mejorar el aprendizaje y aumentar la motivación de los estudiantes. Asimismo, 
\textcite{Rizwan2019} sostienen la necesidad de desarrollar entornos de 
aprendizaje electrónico interactivos para que el alumnado supere los desafíos 
asociados a aprender a programar. De manera similar, \textcite{Hang2023} 
destacan que la RA favorece el aprendizaje, fortaleciendo el pensamiento 
computacional y reduciendo la carga cognitiva, lo que coincide con los 
hallazgos de \textcite{Abidin2020}, quienes la resaltan como una tecnología 
adecuada para captar el interés estudiantil y mitigar la monotonía durante el 
proceso educativo. 

Desde una perspectiva científica, este trabajo de grado aporta a la 
literatura al abordar un vacío de conocimiento y ampliar el marco teórico 
sobre la interacción humano-computadora aplicada a la educación en el ámbito 
nacional. Asimismo, la investigación contribuye a definir los parámetros de 
diseño necesarios para representar conceptos abstractos de la POO en entornos 
tridimensionales, tanto para modelos de RA con marcador como sin marcador, 
lo que permite analizar la influencia de cada variante en el proceso de 
aprendizaje. Con ello, se establece una base teórico-práctica útil para 
futuras investigaciones y para el desarrollo de software inmersivo enfocado 
en este ámbito. 

\section{Alcances y límites}
El presente trabajo de grado abarca el diseño, desarrollo y evaluación de dos 
prototipos funcionales de Realidad Aumentada, uno basado en marcadores físicos 
y otro en posicionamiento espacial sin marcador. Ambos están dirigidos a 
reforzar la asimilación de conceptos de la Programación Orientada a Objetos, 
sin cubrir la totalidad de nociones del paradigma. En cuanto a su alcance 
geográfico, social y temporal, el estudio se delimita a las carreras de 
Ingeniería de Sistemas e Ingeniería Informática de la Facultad de Ciencias y 
Tecnología de la Universidad Mayor de San Simón, mediante la aplicación de 
cuestionarios tipo pre-test y post-test a una muestra de estudiantes de la 
asignatura Introducción a la Programación durante el periodo del primer 
parcial. Cabe aclarar que no se incluye un seguimiento periódico del 
rendimiento académico formal en lo que respecta a calificaciones, ni se busca 
la generalización de los resultados a otras áreas del conocimiento fuera de 
la POO.

\section{Método de investigación}
Pendiente de momento